% Provenance: RESULTADO_RECUPERADO from chapter 24 of the 2249-page integral
% and FORMA_ELIPTICA_RADIO_MODULAR_ACCION_ORIENTADA.md (D04).
% Expository delta: the diagonal operator domain written out fully and
% integer recovery in the two quotient charts made explicit.
% External label dependency: sec:ellipse. Bibliography: elipse,
% hmtedicionintegral, flm, borcherds. No new macros required.

\section{From the action ellipse to T-duality}
\label{sec:ii-dualidad-t}

The reciprocity of the ellipse provides a shape coordinate for the
exchange between visible value and boundary memory. The antecedent is the
same enriched APP--TRIT--TPK state: APP preserves the sum
and product sheets with their remainders and quotients; TRIT preserves their orientation;
TPK transports sheet, carry, route, and memory. The angular and
action outputs already constructed on the joint discrete structure of the continuum
then determine the ellipse. This section composes its subsequent readers:
elliptic shape, lattice exchange, momentum--winding realization,
and modular evaluation \cite{elipse,hmtedicionintegral}.

The cutoff thus lies after the generation of
\(A,C^*,\hbar\). No prescribed radius or string constant selects
those outputs. Boundary memory is first preserved in the
remainder--quotient pair and then transported through the maps proved
below.

\subsection{The action shape determines the normalized radius}

Return to the chart of section~\ref{sec:ellipse}. Let \(A\ne0\),
\(\hbar>0\), and \(|C^*/A|<1\), with both angles expressed in the same
unit. The semiaxes remain labeled:
\begin{equation}
 \eta_{\rm el}=\operatorname{artanh}(C^*/A),\qquad
 a_S=\sqrt{2\hbar}\,e^{\eta_{\rm el}/2},\qquad
 b_S=\sqrt{2\hbar}\,e^{-\eta_{\rm el}/2}.
 \label{eq:ii-t-elipse}
\end{equation}
The balanced phase coordinates and their semiaxes have units of the square root
of action. The radius below is dimensionless.

\begin{proposition}[Shape radius and labeled reciprocity]
\label{prop:ii-t-radio}
The rectangular chart of the ellipse uniquely determines
\begin{equation}
 R_0=\sqrt{\frac{b_S}{a_S}}
 =e^{-\eta_{\rm el}/2}
 =\left(\frac{A-C^*}{A+C^*}\right)^{1/4}>0,
 \qquad \tau_{\rm el}=iR_0^2.
 \label{eq:ii-t-radio}
\end{equation}
It preserves the signed angular ratio through
\begin{equation}
 \frac{C^*}{A}=\frac{1-R_0^4}{1+R_0^4}.
 \label{eq:ii-t-recuperacion-angular}
\end{equation}
Exchanging the labeled semiaxes transforms
\((\eta_{\rm el},R_0,\tau_{\rm el})\) into
\((-\eta_{\rm el},R_0^{-1},-1/\tau_{\rm el})\), leaving
the scale \(a_Sb_S=2\hbar\) fixed.
\end{proposition}

\begin{proof}
Dividing the semiaxes of~\eqref{eq:ii-t-elipse} gives
\(b_S/a_S=e^{-\eta_{\rm el}}\). The map \(R\mapsto R^2\)
is bijective on \(\mathbb R_{>0}\), so its positive root is
unique. For \(x=C^*/A\in(-1,1)\),
\(2\operatorname{artanh}x=\log((1+x)/(1-x))\); hence
\(R_0^4=(1-x)/(1+x)\). Solving for \(x\) proves
\eqref{eq:ii-t-recuperacion-angular}. Exchanging \(a_S,b_S\)
inverts their ratio and changes the sign of \(\eta_{\rm el}\). Finally,
\(-1/(iR_0^2)=iR_0^{-2}\), while the product of the semiaxes
remains \(2\hbar\).
\end{proof}

The angular ratio is recovered; the two absolute angles and their generating
history continue to reside in the enriched state. The point \(R_0=1\)
corresponds to the circular shape \(C^*/A=0\), not to an identification of
all orientations or all histories.

\subsection{Lattice exchange and equivalence of operators}

\begin{definition}[Stable domain of the dual reading]
\label{def:ii-t-reticular}
On \(\mathscr D_0=\mathbb Z^2\times\mathbb R_{>0}\), define
\begin{equation}
 E_{R_0}(m,w)=\frac{m^2}{R_0^2}+w^2R_0^2,\qquad
 \mathsf T_0(m,w,R_0)=(w,m,R_0^{-1}).
 \label{eq:ii-t-energia}
\end{equation}
The first coordinate prolongs the visible reading, and the second the quotient
reading. The full integer domain permits their exchange without requiring
the quotient to belong to a section of nine representatives.
\end{definition}

For the radius of~\eqref{eq:ii-t-radio}, the matrix of this form is
\begin{equation}
 D(\eta_{\rm el})=
 \begin{pmatrix}e^{\eta_{\rm el}}&0\\0&e^{-\eta_{\rm el}}\end{pmatrix}
 =\frac{1}{2\hbar}
 \begin{pmatrix}a_S^2&0\\0&b_S^2\end{pmatrix}.
 \label{eq:ii-t-matriz}
\end{equation}
Indeed, \(R_0^{-2}=e^{\eta_{\rm el}}\) and
\(R_0^2=e^{-\eta_{\rm el}}\). The ellipse equation
\(Q^2/a_S^2+P^2/b_S^2=1\), multiplied by \(2\hbar\), uses
the inverse matrix \(D(-\eta_{\rm el})\). This distinguishes
the squared-semiaxis matrix from the matrix of the elliptic equation.

\begin{theorem}[Lattice duality and unitary conjugation]
\label{thm:ii-t-unitaria}
The transformation \(\mathsf T_0\) is involutive and satisfies
\begin{equation}
 E_{R_0^{-1}}(w,m)=E_{R_0}(m,w).
 \label{eq:ii-t-invariancia}
\end{equation}
On \(\mathcal H_{\rm lat}=\ell^2(\mathbb Z^2)\), let \(H(R_0)\)
be the diagonal operator with values \(E_{R_0}(m,w)\), with domain
\[
 \mathcal D(H(R_0))=
 \left\{c\in\ell^2(\mathbb Z^2):
 \sum_{m,w\in\mathbb Z}E_{R_0}(m,w)^2|c_{m,w}|^2<\infty\right\}.
\]
Then \(H(R_0)\) is self-adjoint and nonnegative. The exchange
\(U_T|m,w\rangle=|w,m\rangle\) is a unitary involution that
transports the domains and satisfies
\begin{equation}
 U_T H(R_0)U_T^{-1}=H(R_0^{-1}).
 \label{eq:ii-t-conjugacion}
\end{equation}
\end{theorem}

\begin{proof}
Exchanging the integers twice and inverting the radius twice returns
the initial point. Moreover,
\[
 E_{R_0^{-1}}(w,m)=w^2R_0^2+m^2R_0^{-2}=E_{R_0}(m,w).
\]
The diagonal domain contains the finitely supported sequences and is dense.
Multiplication by a real sequence is symmetric on that domain.
If \(c\) belongs to the adjoint domain, evaluation against each
basis vector forces the adjoint to have coordinates
\(E_{R_0}(m,w)c_{m,w}\). Their membership in \(\ell^2\) is equivalent
exactly to the condition defining \(\mathcal D(H(R_0))\).
Hence the operator and its adjoint agree. Its nonnegative diagonal values
prove positivity.

The exchange bijectively permutes the orthonormal basis, so it is
unitary and its own inverse. For every sequence \(c\),
\[
 \sum_{m,w}E_{R_0^{-1}}(m,w)^2|(U_Tc)_{m,w}|^2
 =\sum_{m,w}E_{R_0}(m,w)^2|c_{m,w}|^2.
\]
This equality proves exact transport of domains. Finally,
\((U_TH(R_0)U_T^{-1}c)_{m,w}=E_{R_0}(w,m)c_{m,w}
=E_{R_0^{-1}}(m,w)c_{m,w}\); this proves
\eqref{eq:ii-t-conjugacion} on the entire domain, not only on the basis.
\end{proof}

Orientation adds information that energy does not distinguish. In the phase
plane, the exchange \(S_2(Q,P)=(P,Q)\) and rotation
\(J_2(Q,P)=(-P,Q)\) transport the same quadratic form, but
\(S_2^*(dQ\wedge dP)=-dQ\wedge dP\) and
\(J_2^*(dQ\wedge dP)=dQ\wedge dP\). With
\(\lambda=P\,dQ\), one obtains directly
\[
 S_2^*\lambda=d(PQ)-\lambda,\qquad
 J_2^*\lambda=\lambda-d(PQ).
\]
Integration over a closed curve gives opposite actions under \(S_2\)
and equal actions under \(J_2\). Thus the index exchange in
\eqref{eq:ii-t-energia} and the symplectic transport of a trajectory
are different maps, even when they yield the same energy.

\subsection{Nonadic sections and carry transport}

For \(n\in\mathbb Z\), the positive section writes
\[
 r_9^+(n)=1+((n-1)\bmod9),\qquad
 Q_9^+(n)=\frac{n-r_9^+(n)}9,\qquad n=r_9^+(n)+9Q_9^+(n).
\]
The remainder of \(n-1\) is taken in \(\{0,\ldots,8\}\). The change to
the remainder section \(0,\ldots,8\) is
\begin{equation}
 C(r,Q)=\bigl(r-9\varepsilon,Q+\varepsilon\bigr),\qquad
 \varepsilon=\begin{cases}1,&r=9,\\0,&r\ne9.\end{cases}
 \label{eq:ii-t-acarreo}
\end{equation}
Thus \(r+9Q\) remains fixed. For \(n=9\), the two coordinates
are \((9,0)\) and \((0,1)\): their uncorrected energies are
\(81/R_0^2\) and \(R_0^2\). Equality of the integer does not imply
equality of those energy insertions. The following covariant form
resolves the change of section without removing carry.

\begin{definition}[Two quotient charts of the nonadic reading]
\label{def:ii-t-cociente}
Let \(L_9=\mathbb Z(9,-1)\), \(S(m,w)=(w,m)\), and
\(L_9^\vee=SL_9=\mathbb Z(-1,9)\). For
\(L\in\{L_9,L_9^\vee\}\), define
\[
 \mathcal E_{R_0}^{L}([x]_L)
 =\min_{\ell\in L}E_{R_0}(x+\ell).
\]
Transport between charts is
\[
 \mathsf T_{\rm cov}:
 (\mathbb Z^2/L_9)\times\mathbb R_{>0}
 \longrightarrow
 (\mathbb Z^2/L_9^\vee)\times\mathbb R_{>0},\qquad
 ([x]_{L_9},R_0)\longmapsto([Sx]_{L_9^\vee},R_0^{-1}).
\]
\end{definition}

\begin{theorem}[Covariance with crossing memory]
\label{thm:ii-t-cociente}
The preceding energies and transport are well defined. The change
of section~\eqref{eq:ii-t-acarreo} preserves the class in
\(\mathbb Z^2/L_9\), and
\begin{equation}
 \mathcal E_{R_0^{-1}}^{L_9^\vee}([Sx]_{L_9^\vee})
 =\mathcal E_{R_0}^{L_9}([x]_{L_9}).
 \label{eq:ii-t-covariancia}
\end{equation}
The return exchange recovers the initial class and radius. The
integer coordinate maps \(N([x])=x_1+9x_2\) and
\(N^\vee([y])=9y_1+y_2\) are isomorphisms with \(\mathbb Z\)
and satisfy \(N^\vee([Sx])=N([x])\).
\end{theorem}

\begin{proof}
Replacing \(x\) with \(x+\ell_0\), where \(\ell_0\in L\),
permutes the candidates for the minimum. This minimum exists: for \(L=L_9\),
\[
 E_{R_0}(x+k(9,-1))
 =\left(\frac{81}{R_0^2}+R_0^2\right)k^2
  +2\left(\frac{9x_1}{R_0^2}-x_2R_0^2\right)k+E_{R_0}(x),
\]
which tends to infinity with \(|k|\); for \(L_9^\vee\), the
leading coefficient is \(R_0^{-2}+81R_0^2>0\). The change
\(C(x)-x=-\varepsilon(9,-1)\) belongs to \(L_9\).
If \(x'-x\in L_9\), then \(Sx'-Sx\in L_9^\vee\),
proving that transport is well defined. By
\eqref{eq:ii-t-invariancia},
\[
 \min_{\ell\in L_9}E_{R_0^{-1}}(Sx+S\ell)
 =\min_{\ell\in L_9}E_{R_0}(x+\ell).
\]
This proves~\eqref{eq:ii-t-covariancia}; \(S^2=I\) proves
the return between the two charts.

Finally, the homomorphism \(x\mapsto x_1+9x_2\) is
surjective and its kernel is exactly \(L_9\), since
\(x_1=-9x_2\). Similarly, the kernel of
\(y\mapsto9y_1+y_2\) is \(L_9^\vee\). Their induced maps
on the quotients are therefore isomorphisms. Substituting
\(y=Sx\) proves the final identity.
\end{proof}

The recovered integer determines its remainders and quotients once the
section is fixed. The energy-minimizing representative need not be that
canonical representative. The quotient is a covariant representation of the
nonadic reading; the source state also preserves sheet, route, orientation,
boundary, memory, and survival. Nor is \(L_9^\vee\) replaced
by \(L_9\): the exchange is between two transported charts.

\subsection{Momentum, winding, and realization scale}

The circular realization introduces an explicit dimensional chart.
Denote its radius by \(R>0\) and let \(\alpha'>0\) be a parameter
with dimension length squared. This symbol differs from the dimensionless
output \(\alpha_{\rm HMT}\). Fixing the chart does not alter the
angular generator or the shape radius. The normalization and its inverse are
\[
 \mathcal N_{\alpha'}(m,w,R)
 =\left(m,w,\frac{R}{\sqrt{\alpha'}}\right),\qquad
 R=\sqrt{\alpha'}\,R_0.
\]
The indices are now read as momentum number and winding
number; the zero-mode components, expressed in units of
inverse length, are
\begin{equation}
 p_L=\frac mR+\frac{wR}{\alpha'},\qquad
 p_R=\frac mR-\frac{wR}{\alpha'},\qquad
 E_{\rm str}=\frac{p_L^2+p_R^2}{2}
 =\frac{m^2}{R^2}+\frac{w^2R^2}{\alpha'^2}.
 \label{eq:ii-t-quirales}
\end{equation}

\begin{theorem}[Dimensional conjugation of T-duality]
\label{thm:ii-t-cuerda}
The transformation
\[
 \mathsf T_{\alpha'}(m,w,R)=\left(w,m,\frac{\alpha'}R\right)
\]
is an involution and satisfies
\begin{equation}
 \mathcal N_{\alpha'}\mathsf T_{\alpha'}
 =\mathsf T_0\mathcal N_{\alpha'},\qquad
 E_{\rm str}=\alpha'^{-1}E_0\circ\mathcal N_{\alpha'},
 \label{eq:ii-t-normalizacion}
\end{equation}
where \(E_0(m,w,R_0)=E_{R_0}(m,w)\). It preserves \(E_{\rm str}\),
fixes \(p_L\), and reverses the sign of \(p_R\). For fixed
oscillator data \(N_L,N_R\) and intercepts \(a_L,a_R\),
it also preserves the condition
\begin{equation}
 \frac{\alpha'}4(p_L^2-p_R^2)+N_L-N_R-a_L+a_R=0,
 \quad\text{that is}\quad
 N_L-N_R=a_L-a_R-mw.
 \label{eq:ii-t-nivel}
\end{equation}
\end{theorem}

\begin{proof}
Two applications exchange the indices twice and send
\(R\) to \(\alpha'/(\alpha'/R)=R\). Also,
\[
 \mathcal N_{\alpha'}\mathsf T_{\alpha'}(m,w,R)
 =\left(w,m,\frac{\sqrt{\alpha'}}R\right)
 =\mathsf T_0\mathcal N_{\alpha'}(m,w,R).
\]
Substitution into the normalized form gives
\[
 E_0\left(m,w,\frac R{\sqrt{\alpha'}}\right)
 =\frac{\alpha'm^2}{R^2}+\frac{w^2R^2}{\alpha'}
 =\alpha'E_{\rm str}(m,w,R).
\]
For \(R'=\alpha'/R\), the transformed modes are
\[
 p'_L=\frac{wR}{\alpha'}+\frac mR=p_L,\qquad
 p'_R=\frac{wR}{\alpha'}-\frac mR=-p_R.
\]
Thus both their sum and difference of squares are preserved.
Direct expansion gives \(p_L^2-p_R^2=4mw/\alpha'\), transforming
the first equation of~\eqref{eq:ii-t-nivel} into the second.
The exchange preserves \(mw\); with the other data fixed, it preserves
the full condition.
\end{proof}

With these orientations, the abbreviated formula \(N_L-N_R=nw\) requires
\(a_L=a_R\) and \(n=-m\). Retaining the sign in
\eqref{eq:ii-t-nivel} preserves the dictionary between the two readings.
The proved result covers conjugation of radii, zero modes,
energy, and the declared level compatibility. The dimensional interpretation
identifies coordinates through its chart; it does not turn
\(R_0\) alone into a length or infer a complete field theory
from an equality of energies.

\subsection{Modular evaluation and the Moonshine character}

The exceptional construction already developed preserves the incidence of the same
antecedent up to \(V^\natural\). Its graded character has the expression
\(J(\tau)=j(\tau)-744\), with the modular invariance used in
Moonshine \cite{flm,borcherds}. This \(J\) is a function on the
upper half-plane \(\mathbb H\); it is not the rotation \(J_2\), the
exchange \(S\), or the nonadic holonomy \(\Gamma_9\).

\begin{proposition}[Modular semiconjugacy of the radius]
\label{prop:ii-t-modular}
Let \(\iota:\mathbb R_{>0}\to\mathbb H\),
\(\iota(R_0)=iR_0^2\), and
\(S_{\rm mod}(\tau)=-1/\tau\). Then
\begin{equation}
 S_{\rm mod}\circ\iota(R_0)=\iota(R_0^{-1}),\qquad
 J(iR_0^2)=J(iR_0^{-2}).
 \label{eq:ii-t-modular}
\end{equation}
In particular, the two labeled shapes of the ellipse have the same
value under \(J\).
\end{proposition}

\begin{proof}
The first identity is
\(-1/(iR_0^2)=iR_0^{-2}\). The matrix
\(\left(\begin{smallmatrix}0&-1\\1&0\end{smallmatrix}\right)\)
belongs to \(\mathrm{SL}_2(\mathbb Z)\) and acts on
\(\mathbb H\) as \(S_{\rm mod}\). The modular invariance of
\(J\), established in its exceptional construction, gives
\(J(S_{\rm mod}\tau)=J(\tau)\). Applying it to
\(\tau=\iota(R_0)\) yields the second identity.
Proposition~\ref{prop:ii-t-radio} identifies these two arguments
with the moduli of the ellipses of opposite shapes.
\end{proof}

The composition thus preserves an explicit chain:
\[
 (A,C^*,\hbar)\longmapsto(a_S,b_S)
 \longmapsto R_0\longmapsto iR_0^2\longmapsto J(iR_0^2).
\]
Shape inversion is transported to radius inversion and to
\(S_{\rm mod}\); oriented action and change of section retain
their own maps. The function \(J\) assigns a value to a modular class, while
the marking of the axes distinguishes its two presentations and the TPK state
preserves its genealogy. This is the proved relation between the action
ellipse, T-duality, and Moonshine: a composition with domains and
transport identities, whose modular value replaces neither the lattice
operator nor the construction of \(V^\natural\).
